\section*{TÓM TẮT ĐỒ ÁN}
\phantomsection\addcontentsline{toc}{section}{\numberline {}TÓM TẮT ĐỒ ÁN}

Trong bối cảnh ngành công nghiệp làm đẹp và thẩm mỹ phát triển mạnh mẽ ngày nay, việc thấu hiểu và quản lý trải nghiệm khách hàng tại các cơ sở làm đẹp, thẩm mỹ viện đang trở thành một yếu tố cạnh tranh cốt lõi. Tuy nhiên, một thách thức thực tế lớn phát sinh từ việc các phương pháp thẩm mỹ hiện đại (như phẫu thuật thẩm mỹ, tiêm filler, botox, căng chỉ...) có thể làm thay đổi đáng kể cấu trúc và đường nét khuôn mặt của khách hàng sau khi thực hiện liệu trình. Sự thay đổi diện mạo này dẫn đến những khó khăn nghiêm trọng đối với các hệ thống xác thực, nhận dạng và quản lý khách hàng truyền thống khi họ quay lại tái khám hoặc thực hiện các gói dịch vụ tiếp theo.

Đồ án tốt nghiệp này nghiên cứu và phát triển giải pháp \textbf{``Phân tích khách hàng dựa trên biểu cảm khuôn mặt'' (Customer Analysis based on Facial Expressions)} ứng dụng riêng cho lĩnh vực thẩm mỹ và làm đẹp nhằm giải quyết bài toán trên. Bằng việc kết hợp các kỹ thuật thị giác máy tính, học sâu để phân tích điểm mốc (landmarks), biểu cảm và đặc trưng khuôn mặt động, hệ thống có khả năng nhận dạng và xác thực khách hàng một cách chính xác ngay cả khi khuôn mặt có sự thay đổi cơ học do can thiệp thẩm mỹ. Đồng thời, giải pháp tích hợp mô-đun quản lý khách hàng thông minh và bộ trích xuất khuyến nghị tư vấn chăm sóc cá nhân hóa tự động khi khách hàng đến cơ sở, dựa trên đánh giá sinh học về độ hồi phục cơ mặt và mức độ hài lòng. Kết quả nghiên cứu cung cấp một công cụ hỗ trợ đắc lực giúp các thẩm mỹ viện nâng cao chất lượng dịch vụ, nâng cao hiệu quả quản trị và tối ưu hóa hiệu quả kinh doanh.

\newpage
\cleardoublepage